\subsection{CONVOLUTION OF POINT DISTRIBUTIONS AND FUNCTIONS}

Let G be a Lie group, X a manifold of class \(C^{r}\) and \((g, x) \mapsto gx\) a law of left operation of class \(C^{r}\) of G on X. For all \(x \in X\) , let \(\rho(x)\) denote the orbital mapping of x.

DEFINITION 3. Let \(t \in \mathcal{T}^{(s)}(\mathrm{G})\) with \(s \leqslant r\) . Let \(f: \mathrm{X} \to \mathrm{F}\) be a function of class \(\mathbf{C}^r\) with values in a Hausdorff polynormed space (for example \(\mathrm{F} = \mathrm{K}\) ). The convolution of \(t\) and \(f\) , denoted by \(t * f\) , is the function on \(\mathbf{X}\) with values in \(\mathrm{F}\) defined by

\[
(t * f) (x) = \langle t ^ {\vee} * \varepsilon_ {x}, f \rangle .
\]

Then

\[
\begin{array}{l l} (t * f) (x) = \langle \rho (x) _ {*} (t ^ {\vee}), f \rangle & (\text {no. 3, Proposition 14 (ii)}) \\ = \langle t ^ {\vee}, f \circ \rho (x) \rangle & (Diff. \& Anal. Man., R, 13.2.3) \\ = \langle t, (f \circ \rho (x)) ^ {\vee} \rangle & (Diff. \& Anal. Man., R, 13.2.3). \end{array}\tag{4}
\]

Note that Definition 3 can also be written in the more symmetric form

\[
\langle \varepsilon_ {x}, t * f \rangle = \langle t ^ {\vee} * \varepsilon_ {x}, f \rangle .\tag{5}
\]

The function \((g,x)\mapsto f(gx) = (f\circ \rho (x))(g)\) on \(\mathbf{G}\times \mathbf{X}\) is of class \(\mathbf{C}^r\) . By Differentiable and Analytic Manifolds, R, 13.4.4, the function \(x\mapsto \langle t^\vee,f\circ \rho (x)\rangle\) is therefore of class \(\mathbf{C}^{r - s}\) if \(s < \infty\) . In other words, if \(s < \infty\) , \(t*f\) is of class \(\mathbf{C}^{r - s}\) .

Clearly \(t * f\) depends linearly on \(t\) and \(f\) .

Formula (4) implies in particular, for \(g \in G\) ,

\[
\left(\varepsilon_ {g} * f\right) (x) = f \left(g ^ {- 1} x\right)\tag{6}
\]

that is

\begin{enumerate}
\def\labelenumi{(\arabic{enumi})}
\setcounter{enumi}{6}
\tightlist
\item
\end{enumerate}

\[
\varepsilon_ {g} * f = \gamma (g) f.
\]

Suppose that \(\mathbf{K} = \mathbf{R}\) or \(\mathbf{C}\) , that \(\mathbf{G}\) and \(\mathbf{X}\) are finite-dimensional and that \(\mathbf{X}\) has a positive measure invariant under \(\mathbf{G}\) . The definition of \(\varepsilon_g * f\) agrees with that of Integration, Chapter VIII, § 4, no. 1 (cf.~formula (2), loc. cit.).

PROPOSITION 17. Let \(t \in \mathcal{T}^{(s)}(\mathrm{G})\) , \(t' \in \mathcal{T}^{(s')}(\mathrm{X})\) and \(f: X \to \mathrm{F}\) a function of class \(\mathbf{C}^{r}\) with \(s + s' \leqslant r\) . Then

\[
\langle t ^ {\prime}, t * f \rangle = \langle t ^ {\vee} * t ^ {\prime}, f \rangle .
\]

\[
\begin{array}{l l} \langle t ^ {\prime}, t * f \rangle = \langle t ^ {\prime}, x \mapsto \langle t, g \mapsto f (g ^ {- 1} x) \rangle \rangle & \text { by   (4) } \\ = \langle t \otimes t ^ {\prime}, (g, x) \mapsto f (g ^ {- 1} x) \rangle & \text {(Diff. \& Anal. Man., R, 13.4.4)} \\ = \langle t ^ {\vee} \otimes t ^ {\prime}, (g, x) \mapsto f (g x) \rangle & \text {(Diff. \& Anal. Man., R, 13.2.3)} \\ = \langle t ^ {\vee} * t ^ {\prime}, f \rangle . \end{array}
\]

PROPOSITION 18. Let \(t \in \mathcal{T}^{(s)}(\mathrm{G})\) , \(t' \in \mathcal{T}^{(s')}(\mathrm{G})\) and \(f: \mathrm{X} \to \mathrm{F}\) a function of class \(\mathbf{C}^r\) , with \(s + s' \leqslant r\) . Then

\[
(t * t ^ {\prime}) * f = t * (t ^ {\prime} * f).
\]

For all \(x \in X\) ,

\[
\begin{array}{l l} \langle \varepsilon_ {x}, (t * t ^ {\prime}) * f \rangle = \langle (t * t ^ {\prime}) ^ {\vee} * \varepsilon_ {x}, f \rangle & \text {by (5)} \\ = \langle t ^ {\prime \vee} * (t ^ {\vee} * \varepsilon_ {x}), f \rangle & \text {(Propositions 2 and 7)} \\ = \langle t ^ {\vee} * \varepsilon_ {x}, t ^ {\prime} * f \rangle & \text {(Proposition 17)} \\ = \langle \varepsilon_ {x}, t * (t ^ {\prime} * f) \rangle & \text {(Proposition 17).} \end{array}
\]

If \(r \geqslant \infty\) , we see that the set of functions of class \(C^\infty\) on \(X\) with values in \(F\) is a left module over the algebra \(\mathcal{T}^{(\infty)}(G)\) .

PROPOSITION 19. Let \(t \in \mathcal{T}^{(s)}(\mathrm{G})\) , with \(s \leqslant r\) . Let \(f\) (resp. \(f'\) ) be a function of class \(\mathrm{C}^r\) on \(\mathbf{X}\) with values in a Hausdorff polynormed space \(\mathrm{F}\) (resp. \(\mathrm{F}'\) ). Let \((u, u') \mapsto uu'\) be a continuous bilinear mapping of \(\mathrm{F} \times \mathrm{F}'\) into a Hausdorff polynormed space \(\mathrm{F}''\) , so that \(ff'\) is a function of class \(\mathrm{C}^r\) on \(\mathbf{X}\) with values in \(\mathrm{F}''\) . Let \(\sum_{i=1}^{n} t_i \otimes t_i'\) be the image of \(t\) in \(\mathcal{T}^{(s)}(\mathrm{G}) \otimes \mathcal{T}^{(s)}(\mathrm{G})\) under the coproduct. Then

\[
t * (f f ^ {\prime}) = \sum_ {i = 1} ^ {n} (t _ {i} * f) (t _ {i} ^ {\prime} * f ^ {\prime}).
\]

Let \(x \in X\) and let \(\rho(x)\) denote the orbital mapping of x. Then

\[
\begin{array}{l l} \langle \varepsilon_ {x}, t * (f f ^ {\prime}) \rangle = \langle t ^ {\vee}, (f f ^ {\prime}) \circ \rho (x) \rangle & \text {by (4)} \\ = \langle t ^ {\vee}, (f \circ \rho (x)) (f ^ {\prime} \circ \rho (x)) \rangle \\ = \sum_ {i = 1} ^ {n} \langle t _ {i} ^ {\vee}, f \circ \rho (x) \rangle \langle t _ {i} ^ {\prime \vee}, f ^ {\prime} \circ \rho (x) \rangle \\ & (\text {Diff. \& Man. Anal., R, 13.5.2}) \\ = \sum_ {i = 1} ^ {n} \langle \varepsilon_ {x}, t _ {i} * f \rangle \langle \varepsilon_ {x}, t _ {i} ^ {\prime} * f ^ {\prime} \rangle & \text {by (4).} \end{array}
\]

Remark 1. Let G be a Lie group, X a manifold of class \(\mathbf{C}^r\) and \((x,g)\mapsto xg\) a law of right operation of class \(\mathbf{C}^r\) of G on X. If \(t\in \mathcal{T}^{(s)}(\mathrm{G})\) with \(s\leqslant r\) and \(f\colon \mathrm{X}\to \mathrm{F}\) is a function of class \(\mathbf{C}^r\) on X, we denote by \(f*t\) the function on X defined by

\[
\begin{array}{r l} \langle \varepsilon_ {x}, f * t \rangle & = \langle \varepsilon_ {x} * t ^ {\vee}, f \rangle \\ & = \langle \rho (x) _ {*} (t ^ {\vee}), f \rangle \\ & = \langle t ^ {\vee}, f \circ \rho (x) \rangle \\ & = \langle t, (f \circ \rho (x)) ^ {\vee} \rangle . \end{array}\tag{8}
\]

In particular

\[
(f * \varepsilon_ {g}) (x) = f (x g ^ {- 1})\tag{9}
\]

that is

\[
f * \varepsilon_ {g} = \delta (g) ^ {- 1} f.\tag{10}
\]

Propositions 17, 18, 19 become, in the obvious notation,

\begin{enumerate}
\def\labelenumi{(\arabic{enumi})}
\setcounter{enumi}{10}
\tightlist
\item
\end{enumerate}

\[
\langle t ^ {\prime}, f * t \rangle = \langle t ^ {\prime} * t ^ {\vee}, f \rangle\tag{11}
\]

\[
f * (t * t ^ {\prime}) = (f * t) * t ^ {\prime}\tag{12}
\]

\[
(f f ^ {\prime}) * t = \sum_ {i = 1} ^ {n} (f * t _ {i}) (f ^ {\prime} * t _ {i} ^ {\prime}).\tag{13}
\]

PROPOSITION 20. Let \(G\) , \(G'\) be Lie groups, \(X\) a manifold of class \(C^r\) and \((g, x) \mapsto gx\) (resp. \((x, g') \mapsto xg'\) ) a law of left (resp. right) operation of class \(C^r\) of \(G\) (resp. \(G'\) ) on \(X\) . Suppose that \((gx)g' = g(xg')\) for all \(x \in X\) , \(g \in G\) , \(g' \in G'\) . Let \(t \in \mathcal{T}^{(s)}(G)\) , \(t' \in \mathcal{T}^{(s')}(\mathrm{G}')\) and \(f: X \to F\) be a function of class \(C^r\) such that \(s + s' \leqslant r\) . Then

\[
(t * f) * t ^ {\prime} = t * (f * t ^ {\prime}).
\]

For all \(x \in X\) ,

\[
\begin{array}{l l} \langle \varepsilon_ {x}, (t * f) * t ^ {\prime} \rangle = \langle \varepsilon_ {x} * t ^ {\prime \vee}, t * f \rangle & \text { by   (8) } \\ = \langle t ^ {\vee} * (\varepsilon_ {x} * t ^ {\prime \vee}), f \rangle & \text {(Proposition 17)} \\ = \langle t ^ {\vee} * \varepsilon_ {x}, f * t ^ {\prime} \rangle & \text {(Proposition 2 and (11))} \\ = \langle \varepsilon_ {x}, t * (f * t ^ {\prime}) \rangle & \text { by   (5). } \end{array}
\]

In particular, consider G as operating on itself by left and right translations. If \(f \colon \mathbf{G} \to \mathbf{F}\) is a function of class \(\mathbf{C}^r\) on G and \(t \in \mathcal{T}^{(s)}(\mathbf{G})\) (with \(s \leqslant r\) ), \(t * f\) and \(f * t\) are, if \(s < \infty\) , functions of class \(\mathbf{C}^{r - s}\) on G. Further, let \(t' \in \mathcal{T}^{(s')}(\mathbf{G})\) , with \(s + s' \leqslant r\) . Then

\[
(t * f) * t ^ {\prime} = t * (f * t ^ {\prime}).\tag{14}
\]

In particular, \(\mathcal{C}^{\infty}(\mathrm{G})\) is a \((\mathcal{T}^{(\infty)}(\mathrm{G}),\mathcal{T}^{(\infty)}(\mathrm{G}))\) -bimodule. Formulae (5) and (8) admit as special cases

\[
\langle t, f \rangle = \left\langle \varepsilon_ {e}, t ^ {\vee} * f \right\rangle = \left\langle \varepsilon_ {e}, f * t ^ {\vee} \right\rangle .\tag{15}
\]

Remark 2. Let \((g,x)\mapsto gx\) be a law of left operation of class \(\mathbf{C}^r\) of \(G\) on \(X\) . Let \(t\in \mathbf{U}_s(\mathbf{G})\) with \(s\leqslant r\) , \(\Omega\) be an open subset of \(X\) and \(f\colon \Omega \to F\) a function of class \(\mathbf{C}^r\) . \(t*f\) can also be defined by formula (4) or (5); it is a function defined on \(\Omega\) with values in \(F\) , of class \(\mathbf{C}^{r - s}\) if \(s < \infty\) . The results of this no. extend in an obvious way to this situation.

